%==============================================================
%  Masoud (Mas) Hakimi Heris — CV
%  Marginal-label layout (cf. D. Xu, S. Moore)
%  pdflatex cv.tex
%==============================================================
\documentclass[10.5pt,letterpaper]{article}

\usepackage[top=0.7in,bottom=0.8in,left=1.85in,right=0.7in]{geometry}
\usepackage[T1]{fontenc}
\usepackage{lmodern}
\usepackage{enumitem}
\usepackage[hidelinks]{hyperref}
\usepackage{xcolor}
\usepackage{needspace}
\usepackage{microtype}
\usepackage{lastpage}
\usepackage{fancyhdr}

\setlength{\parindent}{0pt}
\setlength{\parskip}{4pt}
\frenchspacing

\pagestyle{fancy}
\fancyhf{}
\renewcommand{\headrulewidth}{0pt}
\fancyfoot[R]{\footnotesize\thepage\ of \pageref{LastPage}}

\setlist[itemize]{leftmargin=1.15em, itemsep=2pt, topsep=3pt, parsep=0pt}

%--------------------------------------------------------------
%  Marginal section label. The label hangs into the left margin;
%  body text continues at the normal text block.
%--------------------------------------------------------------
\newcommand{\cvsec}[1]{%
  \par\needspace{4\baselineskip}\vspace{11pt}%
  \noindent\makebox[0pt][r]{%
    \smash{\parbox[t]{1.45in}{\raggedright\footnotesize\bfseries\MakeUppercase{#1}}}%
    \hspace{0.30in}}%
  \ignorespaces}

% Sub-label inside a section (e.g. "Journal Articles")
\newcommand{\cvsub}[1]{\par\vspace{7pt}\textbf{\itshape #1}\par\vspace{2pt}}

% \entry{left}{right-aligned date}
\newcommand{\entry}[2]{#1\hfill #2\par}

% Numbered publication with hanging indent
\newcommand{\pub}[2]{%
  \par\vspace{4pt}\noindent
  \begingroup\hangindent=2.4em\hangafter=1
  \makebox[2.4em][l]{[#1]}#2\par\endgroup}

\newcommand{\me}{\textbf{M. Hakimi Heris}}
\newcommand{\eqc}{\textsuperscript{*}}
\newcommand{\todo}[1]{{\color{gray}\itshape #1}}

\begin{document}

%==============================================================
%  HEADER
%==============================================================
\noindent\makebox[0pt][r]{\smash{\parbox[t]{1.45in}{\raggedright\bfseries\large
Masoud (Mas)\\ Hakimi Heris}}\hspace{0.30in}}%
{\large\bfseries Ph.D. Candidate, Electrical and Computer Engineering}\\[2pt]
North Carolina State University, Raleigh, NC

\cvsec{Contact\\Information}
\href{mailto:mhakimi@ncsu.edu}{mhakimi@ncsu.edu} \quad +1 (406) 920-9799\\
Web: \href{https://mhakimih1988.github.io/}{mhakimih1988.github.io} \quad
Scholar: \href{https://scholar.google.com/citations?user=VVedtiAAAAAJ}{[link]}\\
ORCID: \href{https://orcid.org/0000-0001-7388-5856}{0000-0001-7388-5856} \quad
GitHub: \href{https://github.com/mhakimih1988}{mhakimih1988} \quad
LinkedIn: \href{https://www.linkedin.com/in/mhakimih/}{mhakimih}\\[2pt]
\todo{Updated: August 2026}

%==============================================================
\cvsec{Research\\Interests}
My research builds the layer between quantum algorithms and the devices that run them:
hardware-native instruction sets and the universality results that justify them, the
compilers and control pulses that target them, and the trapped-ion and photonic hardware
they execute on. I work on quantum algorithms and QSP/QSVT, quantum error analysis and
fault tolerance, hardware--algorithm co-design, adaptive quantum control, hybrid CV--DV
computation, photonic quantum computing, cavity QED, and spin--photon interfaces. My aim
is to make early fault-tolerant quantum computation practical by co-designing the
algorithm, the architecture, and the control together rather than in isolation.

My work falls into four threads, with a representative result in each:

\begin{itemize}
  \item \textbf{Instruction set architectures for hybrid quantum computation}
  [P2]\\
  A native trapped-ion qudit--oscillator ISA whose three core primitives generate
  universal Hamiltonians, with no synthesized mode-coupling operation required; validated
  on dissipative vibronic dynamics compiled onto a single ququart.

  \item \textbf{Pulse-level control and compilation for trapped ions}
  [C4][U2][U3][P1]\\
  \textit{HyPulse}: a GRAPE-based pulse-synthesis framework for hybrid CV--DV
  qubit--oscillator gates, with a reusable library of compound pulse gadgets for
  Hamiltonian simulation.

  \item \textbf{Quantum signal processing: robustness and photonic realization}
  [J1][P3][P4]\\
  Showed that block-encoding and phase-angle errors cross-compensate to first order, and
  built a constrained optimizer suppressing correlated control error by 4--19$\times$;
  co-led the first experimental realization of QSP on a programmable photonic circuit.

  \item \textbf{Nanophotonics and quantum materials}
  [C2][C3][J2][J3]\\
  Inverse-designed nanophotonic cavities for spin--photon interfacing, and MEMS/NEMS
  platforms for optical strain studies of 2D materials.
\end{itemize}

%==============================================================
\cvsec{Education}
\entry{\textbf{Ph.D. Candidate, Electrical Engineering}, North Carolina State University}{2024--Present}
Advisor: Prof.\ Yuan Liu \quad$\cdot$\quad Expected May 2027 \quad$\cdot$\quad GPA: 4.0/4.0\\
Dissertation: \emph{Co-Design of Quantum Algorithms, Architectures, and Control for
Early Fault-Tolerant Hardware}

\vspace{5pt}
\entry{\textbf{M.S., Electrical Engineering}, Montana State University}{2022--2024}
Advisor: Prof.\ David L. Dickensheets \quad$\cdot$\quad GPA: 3.87/4.0\\
Thesis: \emph{MEMS Cantilevers for Dynamic Strain Studies of 2D Materials}

\vspace{5pt}
\entry{\textbf{M.S., Electrical Engineering (Electronics)}, Qazvin Azad University, Iran}{2012--2015}
\entry{\textbf{B.S., Electrical Engineering (Electronics)}, Karaj Azad University, Iran}{2006--2011}

%==============================================================
\cvsec{Publication\\Summary}
\textbf{Published:} 7 peer-reviewed papers, 2 of them first-authored.
\textbf{In progress:} 3 preprints or manuscripts under review and 4 in preparation,
5 of these 7 first- or co-first-authored.\\
\textbf{Impact:} 90 citations, h-index 3, i10-index 2 (Google Scholar, August 2026).\\
Venues include IEEE QCE, HPCA (under review), ISCA workshops, SPIE Photonics for Quantum,
and \emph{APL Computational Physics}.

%==============================================================
\cvsec{Peer-Reviewed\\Publications}
\cvsub{Journal Articles}

\pub{J3}{E. Karooby, H. Sahbafar, \me, A. Hadi, V. Eskandari.
``Identification of Low Concentrations of Flucytosine Drug Using a SERS-Active Filter
Paper Substrate.'' \emph{Plasmonics} 19(2), 855--863, 2024.}

\pub{J2}{V. Eskandari, H. Sahbafar, E. Karooby, \me, S. Mehmandoust, D. Razmjoue, A. Hadi.
``SERS Filter Paper Substrates Decorated with Silver Nanoparticles for the Detection of
Molecular Vibrations of Acyclovir Drug.''
\emph{Spectrochimica Acta Part A} 298, 122762, 2023.}

\pub{J1}{K. J. Joven, E. R. Das, J. Bierman, A. Majumdar, \me, Y. Liu.
``Scalable Quantum Computational Science: A Perspective from Block-Encodings and
Polynomial Transformations.'' \emph{APL Computational Physics} 2, 010901, 2026.}

\cvsub{Conference Papers}

\pub{C4}{\me, Y. Liu, F. Mueller.
``HyPulse: A Pulse Synthesis Framework for Hybrid Qubit--Oscillator Gates on Trapped-Ion
Platform.'' \emph{IEEE International Conference on Quantum Computing and Engineering
(QCE)}, 2026. arXiv:2604.26804.}

\pub{C3}{E. Karooby, \me, Q. Gu, D. Farfurnik.
``Inverse Design of Non-Periodic Circular Cavities for Efficient Spin--Photon Interfacing
by Genetic Algorithm.'' \emph{SPIE Optical Modeling and Performance Predictions XV}
13594, 116--119, 2025.}

\pub{C2}{A. Das, N. K. Vij, \me, D. Farfurnik.
``Optimization of Circular Gratings via Guided-Mode-Expansion-Method-Based Inverse
Design.'' \emph{SPIE Photonics for Quantum} 13563, 65--71, 2025.}

\pub{C1}{\me, R. Faez.
``Investigation of the Possibility of Using Single-Electron Transistors in Digital-to-Time
Converters.'' \emph{IEEE Global Congress on Electrical Engineering (GC-ElecEng)},
150--153, 2021.}

%==============================================================
\cvsec{Under Review\\and Preprints}
\pub{U3}{C. Chu, \me, H. Zhou, Y. Liu.
``FACT: A Fidelity-Aware Compiler for Trapped-Ion Quantum Computers.''
Under review, \emph{IEEE International Symposium on High-Performance Computer
Architecture (HPCA)}, 2027.}

\pub{U2}{R. Patel, \me, Y. Liu, F. Mueller.
``Synthesizing Compound Pulse Gadgets for Hamiltonian Simulation on Trapped-Ion
Platforms.'' arXiv:2607.00826, 2026.}

\pub{U1}{\me.
``Graph State Generation Based on Quantum Photonic Devices, Enabling Measurement-Based
Quantum Computing.'' arXiv:2607.16667, 2026.}

%==============================================================
\cvsec{Manuscripts in\\Preparation}
\pub{P4}{E. Karooby\eqc, \me\eqc, Y. Liu, Q. Gu.
``Universal Function Approximation Using Quantum Signal Processing on a Photonic
Integrated Platform.'' \quad \eqc\,Equal contribution.}

\pub{P3}{\me, [co-authors], Y. Liu.
``Robustness-Aware Quantum Signal Processing: Error Cancellation, Optimization, and
Learned Angle Synthesis.''}

\pub{P2}{\me, Y. Liu, F. Mueller.
``A Native Trapped-Ion Hybrid Qudit--Oscillator Instruction Set.''}

\pub{P1}{\me, Y. Liu, F. Mueller.
``Learned Pulse Synthesis for Hybrid CV--DV Gates on Trapped-Ion Hardware.''}

%==============================================================
\cvsec{Funding and\\Projects}
\textbf{Grand Unification of Quantum Algorithms on Photonic Integrated Circuits}
\begin{itemize}
  \item Source of Support: NSF I-Corps Teams (Quantum+X)
  \item Role: Co-Entrepreneurial Lead
  \item Project Dates: 2026
  \item PI: Qing Gu (NC State, ECE/Physics); Technical Co-Lead: Yuan Liu (NC State, ECE/CS);
        Co-Entrepreneurial Lead: Elaheh Karooby; Industry Mentor: Amirhassan Shams-Ansari
        (Monarch Quantum)
  \item Total Project Amount: \$50,000
\end{itemize}

\vspace{7pt}
\textbf{NQVL:QTSD:Pilot/Design: Quantum Advantage-Class Trapped Ion System (QACTI)}
\begin{itemize}
  \item Source of Support: NSF National Quantum Virtual Laboratory, via subcontract from
        Duke University
  \item Role: Supported as Graduate Research Assistant
  \item Lead PI: Ken Brown (Duke University); NC State: Yuan Liu, Frank Mueller
  \item Phase 1 (NSF 2410675): \$100,000, 08/2024--07/2025
  \item Phase 2 (NSF 2531350): \$390,000, 09/2025--08/2027
  \item Building an advantage-class trapped-ion quantum computer for remote use by the
        scientific community, with an open-source hardware/software stack.
\end{itemize}

\vspace{4pt}
\todo{[Also acknowledged in your papers: NSF PHY-2325080, OMA-2120757, and DOE
DE-SC0025384. Add if they support your position.]}

%==============================================================
\cvsec{Selected\\Presentations}
\entry{``Experimental Realization of Quantum Signal Processing on a Photonic Integrated
Circuit,'' IEEE Photonics Conference.}{2026}
\entry{``Toward Compact Multiqubit Gates on Trapped Ions: A Hardware-Aware Pulse
Synthesis,'' QCCC-26 Workshop, ISCA.}{2026}
\entry{``Simulations of Nonadiabatic Vibronic Dynamics on Hybrid Qudit--Qumode Trapped-Ion
Devices,'' RQS Institute Workshop, University of Maryland.}{2026}
\entry{``MEMS Devices for In-Situ Straining of 2D Materials,'' MonArk NSF Quantum Foundry
Site Visit, Bozeman, MT.}{2024}
\entry{``MEMS Cantilevers for Dynamic Strain Studies of 2D Materials,'' Optical Science
\& Engineering Conference (OpTeC), Bozeman, MT.}{2023}

%==============================================================
\cvsec{Awards and\\Professional\\Programs}
\entry{Prospects in Theoretical Physics (PiTP), Institute for Advanced Study, Princeton}{2026}
\entry{Graduate Merit Award, College of Engineering, NC State University}{2024--2025}
\entry{AIM Photonics Summer Academy, MIT}{2023}
\entry{Summer of Quantum Short Course, LPS Qubit Collaboratory}{2023}
\entry{Excellent Graduate Student in Electronics Engineering, Qazvin Azad University}{2015}
\entry{Top 3 Undergraduate Student in Electronics Engineering, Karaj Azad University}{2009}
\entry{Best Patent, Kharazmi Festival, Tehran, Iran}{2004}

%==============================================================
\cvsec{Research\\Experience}
\entry{\textbf{Graduate Researcher}, Quantum Computing and Engineering, NC State University}{2024--Present}
Advisor: Prof.\ Yuan Liu
\begin{itemize}
  \item Designed a native trapped-ion hybrid qudit--oscillator instruction set and proved
  that three of its primitives generate universal Hamiltonians, with no synthesized
  mode-coupling operation required.
  \item Compiled dissipative vibronic dynamics of a three-chromophore complex onto a single
  $^{171}$Yb$^+$ ququart, cutting gate count from 96 to 28 per Trotter step relative to the
  superconducting realization and reproducing published trajectories to ${\sim}1.1\%$ RMSE.
  \item Built HyPulse, a GRAPE-based pulse-synthesis framework for hybrid CV--DV gates, and
  extended HybridLane (PennyLane) to native qudit--oscillator circuits.
  \item Co-led the first experimental realization of QSP on a programmable photonic
  integrated circuit, synthesizing nonlinear functions on linear-optical hardware.
  \item Optical characterization and design of nanophotonic structures for light--matter
  interaction (cavity QED).
  \item Collaborations across ECE, physics, and computer science, including Duke
  University's NQVL for trapped-ion experimental validation.
\end{itemize}

\vspace{5pt}
\entry{\textbf{Graduate Researcher}, MonArk NSF Quantum Foundry, Montana State University}{2022--2024}
Advisors: Prof.\ David L. Dickensheets, Prof.\ Nicholas J. Borys
\begin{itemize}
  \item Designed, fabricated, and characterized MEMS/NEMS devices for optical studies of
  2D quantum materials, from EBL and photolithography through SEM/AFM metrology.
\end{itemize}

%==============================================================
\cvsec{Teaching\\Experience}
\entry{\textbf{Preparing the Professoriate (PTP) Fellow}, NC State University}{2026--2027}
Year-long teaching apprenticeship: co-teach a course with a faculty mentor, develop a teaching
portfolio and philosophy statement, and complete pedagogy coursework.

\vspace{5pt}
\entry{Teaching Assistant, Signals, Circuits, and Systems Laboratory, NC State University}{[term]}
\entry{Teaching Assistant, Microfabrication and Electrical Circuits Laboratories, Montana State University}{2022--2023}

\vspace{4pt}
\todo{[Add enrollments and what you ran. For faculty applications, also list the courses
you are prepared to teach.]}

%==============================================================
\cvsec{Mentoring}
\entry{\textbf{Samuel Firmansyah} --- [level], NC State University}{Spring 2026}
Project: AI-based robust quantum signal processing (learned robust angle synthesis).

\vspace{4pt}
\entry{\textbf{Nidhi Grover} --- [level], NC State University}{Spring 2026}
Project: AI-based pulse synthesis for hybrid CV--DV gates on trapped-ion hardware.

%==============================================================
\cvsec{Professional\\Service}
\cvsub{Peer Review}
Sub-reviewer on program-committee and editorial assignments for Prof.\ Yuan Liu and
Prof.\ Frank Mueller (5 manuscripts, 2025--2026):
\begin{itemize}
  \item \emph{IEEE Journal of Selected Topics in Signal Processing}, Special Issue on
        Quantum Technologies --- 1 manuscript (2026)
  \item IEEE International Conference on Quantum Computing and Engineering (QCE, IEEE
        Quantum Week) --- 2 manuscripts, Quantum Applications (QAPP) and Quantum Machine
        Learning tracks (2026)
  \item International Workshop on Software Frameworks and Workload Management on Quantum
        and HPC Ecosystems (SFWM), at SC --- 2 manuscripts (2025, 2026)
\end{itemize}

\cvsub{Memberships}
\todo{[IEEE, APS, Optica, ACM --- list any current student memberships.]}

\cvsub{Departmental and Outreach}
\todo{[Seminar or reading-group organizing, poster judging, recruiting/visit-day hosting,
graduate student association, K--12 or undergraduate outreach. Anything counts; this
subsection is still empty.]}

%==============================================================
\cvsec{Technical\\Skills}
\textbf{Quantum algorithms \& theory:} QSP/QSVT, block encodings, Hamiltonian simulation,
open quantum systems, hybrid CV--DV computation, quantum error analysis.\\[2pt]
\textbf{Quantum control:} GRAPE, pulse synthesis, constrained optimization, differentiable
simulation, hardware-aware compilation, Hamiltonian-based control.\\[2pt]
\textbf{Software:} Python, JAX, QuTiP, PennyLane, Qiskit, Strawberry Fields, MATLAB,
Julia, C/C++.\\[2pt]
\textbf{Machine learning:} Neural networks, Fourier neural operators, learned pulse
synthesis, physics-informed and differentiable training.\\[2pt]
\textbf{Photonics \& hardware:} Programmable photonic circuits, Lumerical
FDTD/INTERCONNECT/qINTERCONNECT, COMSOL, nanophotonic design and fabrication, optical
characterization.

%==============================================================
\cvsec{Earlier Industry\\Experience}
Software and electrical engineering roles, 2011--2021 (FDM4 Canada; power and industrial
systems engineering and project management).

%==============================================================
% \cvsec{References}
% \textbf{Prof.\ Yuan Liu} --- NC State University --- q\_yuanliu@ncsu.edu --- Ph.D. advisor\\
% \textbf{Prof.\ Frank Mueller} --- NC State University --- fmuelle@ncsu.edu\\
% \textbf{Prof.\ David L. Dickensheets} --- Montana State University --- davidd@montana.edu --- M.S. advisor
%==============================================================

\end{document}
